\documentclass[a4paper]{article}
\usepackage[utf8]{inputenc} % Skal passe til editorens indstillinger
\usepackage[english]{babel} % danske overskrifter


\newcommand{\name}{Søren Krogh Andersen \& Carsten Lau Nielsen}
\newcommand{\stnumber}{s123369, s123161}
\newcommand{\course}{Special Course}
\newcommand{\university}{Danmarks Tekniske Universitet}
\newcommand{\studyline}{DTU Elektroteknologi}
\newcommand{\assignment}{Attitude controller for multirotor}
\renewcommand{\date}{\today} %If another date, than that of today is desiered


% Palatino for rm and math | Helvetica for ss | Courier for tt
\usepackage{mathpazo} % math & rm
\linespread{1.05}        % Palatino needs more leading (space between lines)
\usepackage{palatino} % tt
\normalfont
\usepackage[T1]{fontenc}

\usepackage{graphicx}%allerese hentet % indsættelse af billeder
\usepackage{epstopdf} %Tilfj "--enable-write18" i argumentet for LaTex build. Dette vil konvertere .eps figurer til pdf-format
\graphicspath{{./picture/}} % stivej til bibliotek med figurer
\usepackage{subcaption} %Til gruppering af figurer
\usepackage{amsmath} %matpakke
\usepackage{amsfonts} %
\usepackage{amssymb} %
\usepackage{steinmetz} % flere matematik symboler
\usepackage{polynom} %for displaying polynom division
\usepackage{mathtools} % matematik - understøtter muligheden for at bruge \eqref{}
\usepackage{float}
\usepackage{placeins}
\usepackage{hhline}

%
\usepackage[usenames,dvipsnames]{xcolor}
\usepackage[compact,explicit]{titlesec}% http://ctan.org/pkg/titlesec
%

%---------%
%Easy edit%
%---------%

%Section formating. arg1 is supplied when making section
\newcommand\presectionnumber[1]{~~}
\newcommand\postsectionnumber[1]{}
\newcommand\midlesection[1]{#1}
\newcommand\sectionnum[1]{\arabic{#1}}
\newcommand\subsectionnum[1]{\arabic{#1}}
\newcommand\subsubsectionnum[1]{\alph{#1}}



%------------%
%setion setup%
%------------%
\renewcommand\thesection{Opgave~\sectionnum{section}} %pas p�, kun i matematik
\renewcommand\thesubsection{\thesection,~\subsectionnum{subsection}}
\definecolor{MagRed}{RGB}{190,40,15}
\definecolor{MathGreen}{RGB}{82,164,0}

\titleformat{\section}{\normalfont\sffamily\large\bfseries\color{MathGreen}}{}{0pt}{|\kern-0.15ex|\kern-0.15ex|\kern-0.15ex|\presectionnumber{#1}\sectionnum{section}\postsectionnumber{#1}\qquad\quad\midlesection{#1}\label{sec:\sectionnum{section}}}
\titleformat{\subsection}[runin]{\large\bfseries}{}{10pt}{\sectionnum{section}.\subsectionnum{subsection})~#1\label{sec:\sectionnum{section}.\subsectionnum{subsection}}}
\titleformat{\subsubsection}[runin]{\itshape}{}{0pt}{\subsectionnum{subsection},\subsubsectionnum{subsection}~#1\label{sec:\sectionnum{section}.\subsectionnum{subsection}.\subsubsectionnum{subsubsection}}}
%\titleformat{\subsubsection}{\bfseries}{}{0pt}{\alph{subsection}.\arabic{subsubsection})\qquad\quad#1\label{\arabic{section}\alph{subsection}\arabic{subsubsection}}}

%----------%
%page setup%
%----------%
\textwidth = 400pt
\marginparwidth = 86pt
\hoffset = -25pt
\voffset= -30pt
\textheight = 670pt

%--------%
%hyperref%
%--------%
\newcommand{\HRule}{\rule{\linewidth}{0.5mm}}
\usepackage{fancyhdr}
\usepackage[plainpages=false,pdfpagelabels,pageanchor=false]{hyperref} % aktive links
\hypersetup{%
  pdfauthor={\name - \stnumber},
  pdftitle={\assignment},
  pdfsubject={\course} }
%\usepackage{memhfixc}% rettelser til hyperref

%-------------%
%Headder setup%
%-------------%
\fancyhf{} % tom header/footer
\fancyhfoffset{20pt}
\fancyhfoffset{20pt}
\fancyhead[OL]{\name \\ \course}
\fancyhead[OC]{Date \\ \date}
\fancyhead[OR]{\university\\ \studyline}
\fancyfoot[FL]{}
\fancyfoot[FC]{\thepage}
\fancyfoot[FR]{}
\renewcommand{\headrulewidth}{0.4pt}
\renewcommand{\footrulewidth}{0.4pt}
\pagestyle{fancy}
 % style setup

%Listings%
\usepackage{listingsutf8}
\usepackage[framed,numbered]{matlab-prettifier}


%setup listings
\lstset{language=Matlab,
  extendedchars=true,
  language=Octave,                % the language of the code
  basicstyle=\ttfamily\footnotesize,           % the size of the fonts that are
  % used for the code
  numbers=left,                   % where to put the line-numbers
  numberstyle=\tiny\color{gray},  % the style that is used for the line-numbers
  stepnumber=2,                   % the step between two line-numbers. If it's 1, each line 
                                  % will be numbered
  numbersep=5pt,                  % how far the line-numbers are from the code
  backgroundcolor=\color{white},      % choose the background color. You must add \usepackage{color}
  showspaces=false,               % show spaces adding particular underscores
  showstringspaces=false,         % underline spaces within strings
  showtabs=false,                 % show tabs within strings adding particular underscores
  frame=single,                   % adds a frame around the code
  rulecolor=\color{black},        % if not set, the frame-color may be changed on line-breaks within not-black text (e.g. comments (green here))
  tabsize=4,                      % sets default tabsize to 2 spaces
  captionpos=b,                   % sets the caption-position to bottom
  breaklines=true,                % sets automatic line breaking
  breakatwhitespace=false,        % sets if automatic breaks should only happen at whitespace
  title=\lstname,                   % show the filename of files included with \lstinputlisting;
                                  % also try caption instead of title
  %keywordstyle=\color{blue},          % keyword style
  %commentstyle=\color{dkgreen},       % comment style
  %stringstyle=\color{mauve},         % string literal style
  escapeinside={\%*}{*)},            % if you want to add LaTeX within your code
  morekeywords={*,...},              % if you want to add more keywords to the set
  deletekeywords={...}              % if you want to delete keywords from the given language
}
\lstset{literate=
  {á}{{\'a}}1 {é}{{\'e}}1 {í}{{\'i}}1 {ó}{{\'o}}1 {ú}{{\'u}}1
  {Á}{{\'A}}1 {É}{{\'E}}1 {Í}{{\'I}}1 {Ó}{{\'O}}1 {Ú}{{\'U}}1
  {à}{{\`a}}1 {è}{{\`e}}1 {ì}{{\`i}}1 {ò}{{\`o}}1 {ù}{{\`u}}1
  {À}{{\`A}}1 {È}{{\'E}}1 {Ì}{{\`I}}1 {Ò}{{\`O}}1 {Ù}{{\`U}}1
  {ä}{{\"a}}1 {ë}{{\"e}}1 {ï}{{\"i}}1 {ö}{{\"o}}1 {ü}{{\"u}}1
  {Ä}{{\"A}}1 {Ë}{{\"E}}1 {Ï}{{\"I}}1 {Ö}{{\"O}}1 {Ü}{{\"U}}1
  {â}{{\^a}}1 {ê}{{\^e}}1 {î}{{\^i}}1 {ô}{{\^o}}1 {û}{{\^u}}1
  {Â}{{\^A}}1 {Ê}{{\^E}}1 {Î}{{\^I}}1 {Ô}{{\^O}}1 {Û}{{\^U}}1
  {œ}{{\oe}}1 {Œ}{{\OE}}1 {æ}{{\ae}}1 {Æ}{{\AE}}1 {ß}{{\ss}}1
  {ç}{{\c c}}1 {Ç}{{\c C}}1 {ø}{{\o}}1 {å}{{\r a}}1 {Å}{{\r A}}1
  {€}{{\EUR}}1 {£}{{\pounds}}1
}

 \lstloadlanguages{% Check Dokumentation for further languages ...
         %[Visual]Basic
         %Pascal
         %C
         %C++
         %XML
         %HTML
         %Java
         %VHDL
         Matlab
 }
 %Listings slut%









%Matematik hurtige ting
%fed
\renewcommand\vec[1]{\mathbf{#1}}
\newcommand\matr[3]{{}_{#2}\mathbf{#1}{}_{#3}}
\newcommand\facit[1]{\underline{\underline{#1}}}
%\renewcommand\d[3]{\frac{\mbox{d}^{#3}#1(#2)}{\mbox{d}#2^{#3}}}
%underline
%\renewcommand\vec[1]{\underline{#1}}
%\newcommand\matr[3]{{}_{#2}\underline{\underline{#1}}{}_{#3}}

\renewcommand\matrix[4]{ %{alignment}{to space}{from space}{matrix}
{\vphantom{\left[\begin{array}{#1}#4\end{array}\right]}}_{#2}\kern-0.5ex
\left[\begin{array}{#1}
#4
\end{array}\right]_{#3}
}
\newcommand\e[0]{\mbox{e}}
\newcommand\E[1]{\cdot 10^{#1}}
\newcommand\im[0]{i}

\newcommand\Jaco{\mbox{Jacobi}}
\newcommand\del[2]{\frac{\partial {#1}}{\partial {#2}}}
\newcommand\abs[1]{\left| {#1} \right|}
\newcommand\stdfig[4]{ %width,img,cap,lab
\begin{figure}[H]
\centering
\includegraphics[width={#1}\textwidth]{#2}
\caption{#3}
\label{#4}
\end{figure}
}
\newcommand\stdfignoscale[3]{ %img,cap,lab
\begin{figure}[H]
\centering
\includegraphics{#1}
\caption{#2}
\label{#3}
\end{figure}
}
\newcommand\diff{\dot}
\newcommand\ddiff{\ddot}
\newcommand\dddiff{\dddot}
\newcommand\ddddiff{\ddddot}






% How to make ref to books or urls in bib
%\citetitle[fx: page 1]{name of ref in bib}

\begin{document}
\begin{titlepage}
\centering \parindent=0pt

\vspace*{\stretch{1}} \HRule\\[1cm]\Huge
\course\\[0.7cm] 
\large 
Supervisor: Søren Hansen sh@elektro.dtu.dk
%\assignment\\[1cm]
\HRule\\[4cm]  
%\includegraphics[width=6cm]{picture}\\ Use this if you want a picture on the frontpage
\name\\
\stnumber

\vspace*{\stretch{2}} \normalsize %
\begin{flushleft}

\date \end{flushleft}
\vspace*{\stretch{2}} \normalsize
\begin{flushright}
\includegraphics[width=6cm]{./dtu.eps}\\
\end{flushright}
\end{titlepage}

\newpage

%Quick hardware walktrough
\section{ Introduction }
The aim of this project was to create an attitude controller with a rappid response for a, potential, autonomus drone.
The goal of this project was also to function as a primer, for a possible future project. This should allow for potential difficoulties and obsticles to show early. As it was the case wich issues in form of mechanical vibrations from motors and propellers.
Half of the time was used on constructing and programming the hardware, while the other half was used for designing and implementing an attitude controller and investigating visual odometry methods.
As the investigations in visual odometry did not produce any results, apart from experience, these are not documented in this report.



\section{Hardware description}
Most of the hardware was contructed before the start of the course, mainly
because we needed a custom PCB fabricated and lead times for PCBs from China are
unpredictable and can be long.
Additionally, a lot of 3D printing was needed for the frame, which takes quite
some time (total printing time for the parts used was about 14h, then comes
part design, prototypes and printer setup).

\subsection{ Flight controller }
The main processor of the fligt controller is a \emph{Texas
Instruments} \emph{TM4C123GH6 tiva series ARM CORTEX M4F} processor.
This processor was chosen due to ease of use.
The chip is avaliable from most distributors and cheap development boards,
\emph{TI Launchpad}s, can be used as debuggers. Furthermore a huge driver
library (\emph{TivaWare}) for hardware abstraction is available from \emph{TI}s
webpage.

The flighcontroller has a varity of IO ports: PWM outputs, I2C, CAN, UART,
analog and digital inputs, a special port for conneting to the \emph{Hardkernel
oDroid U3} Linux computer via SPI, SPI connection for a radio module, and
onboard IMU (\emph{Invensense mpu9250}) and barometer (\emph{Freescale
MPL3115A2}).
\begin{figure}
	\centering
	\includegraphics[width=0.7\textwidth]{./picture/flightcontroller.JPG}
	\caption{The populated flight controller PCB.}
	\label{fig:flightcontroller}
\end{figure}

A close-up of the populated flight controller PCB can be seen in figure
\ref{fig:flightcontroller}. Note that not all IO modules are populated with the
needed parts. Some mistakes were made in the layout, and thus had to be fixed
(pink wire nest).

\subsection{ Radio }
We chose the \emph{TI CC1101} $433MHz$ radio module, as this was relativly easy
to get working.

For non-tethered telemetry a remote controller was contructed, capable of
sending low amounts of data between the flightcontroller and remote. The remote
has a USB port for dumping telemetry data to a PC and joystics for sending
steering commands to the flight controller.

A picture of the crude radio remote controller can be seen in figure
\ref{fig:remote}

\begin{figure}
	\centering
	\includegraphics[width=0.7\textwidth]{./picture/remote.JPG}
	\caption{The remote control with analog joysticks and status screen.}
	\label{fig:remote}
\end{figure}

\subsection{ Motors, ESCs, propellers }
The \emph{EMAX MT1806} motors
and recommended drivers for those were chosen due to good reviews from
hobbyists. Likewise 5030 carbonfiber propellers were chosen, as they are the
recomended propellers for motors of this size.

\section{ System identification } %change title?
While the differential equations governing attitude are highly nonlinear,
they can be linearized about an operating point, where the angular speed is zero.
Thus the gyroscopic effects can be omitted.

The thrust of each motor is assumed to be linear with the control output and the
retardation can be modelled as a first order lowpass filter.

The differential equations for rotation around the $x$ and $y$ axes
becomes:
\begin{equation}
 \dddiff\theta_x = \frac{1}{\tau_{motor}} \cdot \left(
 (u_1 + u_2 - u_3 - u_4) G_{thrust} l_{y} \frac{ 1 }{ I_{x} }  - \ddiff\theta_x
 \right) 
\end{equation}
\begin{equation}
 \dddiff\theta_y = \frac{1}{\tau_{motor}} \cdot \left(
 (-u_1 + u_2 + u_3 - u_4) G_{thrust} l_{x} \frac{ 1 }{ I_{y} }  - \ddiff\theta_y
 \right) 
\end{equation}
Where $G_{thrust}$ is the thrust in $N$ produced by a single propeller at full
thrust, $l_{x}$ is the length from the center of mass to the motors along the
$x$-axis, $I_{y}$ is the moment of inertia around the $y$-axis,
$\tau_{motor}$ is the time constant of the motor and $u_{[1;4]}$ are motor
commands for motor $[1;4]$ in the range $[0;1]$.

Local coordinate system for the drone, and motor numbering is chosen as shown in
figure \ref{fig:axis}
\stdfig{0.8}{axis}{Choice of local coordinate system and motor numbering for
drone}{fig:axis}

We control the motors in tandem, so:
\begin{align}
u_1 = & u_x - u_y + T\\
u_2 = & u_x + u_y + T\\
u_3 = & -u_x + u_y + T\\
u_4 = & -u_x - u_y + T\\
\end{align}
Where T is a general thrust command the equations simplify to:
\begin{equation}
 \dddiff\theta_x = \frac{1}{\tau_{motor}} \cdot \left(
 4 u_x G_{thrust} l_{y} \frac{ 1 }{ I_{x} }  - \ddiff\theta_x
 \right) 
\end{equation}
\begin{equation}
 \dddiff\theta_y = \frac{1}{\tau_{motor}} \cdot \left(
4 u_y G_{thrust} l_{x} \frac{ 1 }{ I_{y} }  - \ddiff\theta_y
 \right) 
\end{equation}


\section{ Sensor filtering } %change title?
Since the rotating propellers introduce a lot of vibration some sort of
filtering is required. Clasically mechanical filters (foam and rubber-rings) have
been used as a solutuion to this problem. We propose using a high-samplerate
IMU and using a digital filter to supress vibration above the output rate of the
controller. Assuming that the plant dynamics are temporarily linear, we can
perform further state estimation on a downsampled version of the IMU- readings,
allowing for more complex estimators.

As an attitude estimator we use the popular \emph{Madgwick sensor fusion
algorithm}\cite{Madgwick2011}, which
fuses measuremenst form a 3-axis gyroscope and 3-axis accelerometer. The gyro measuremens are integrated to find the attitude and the
accelerometer is used to compensate for gyro drift.

The angular acceleration can be estimating by nummerical differentiation of the
filtered gyro measurements.

\section{ Controller } %change title?

If the input to the attitude is the desired thrust vector, we can measure the
angular differences as shown in figure \ref{fig:controller}.

We denote this angle $\vec{\Delta\theta} = \matrix{c}{}{}{ \Delta \theta_x \\
\Delta\theta_y}$ and the angular rotation and angular acceleration
in the local coordinate system: $\vec\omega$ and $\vec\alpha$.

\stdfig{0.8}{controller}{Diagram showing how the difference in
current and desired angle is calculated}{fig:controller}

The controller for each axis is designed via pole placement, with state
limitation, as in~\cite{WangWang} 
, but without the refference model.

The poles are placed to form a critically dampened $3^{rd}$ system, with triple
pole in $\omega$, this system can be seen in figure \ref{fig:controller-no-fix}
where:
\begin{align}
k_0 & = \omega^3\\
k_1 & = 3\omega^2\\
k_2 & = 3\omega
\end{align}

\stdfig{0.8}{controller-no-fix}{Diagram for desired
system}{fig:controller-no-fix}

To hide the rotor dynamics, we introduce two new gains: $g_1$ and $g_2$ as seen
in figure \ref{fig:controller-fix}, where $G$ is the static gain of the system
given by max thrust, arm length etc.
This makes the rotor dynamics behave, as if the time constant was $\frac{1}{3\omega}$ instead of $\tau_{motor}$.
The gains $g_1$ and $g_2$ become:
\begin{align}
g_1 = k_2 \tau_{motor,est}
g_2 =1 - \frac{1}{g_1}
\end{align}

\stdfig{0.9}{controller-fix}{Diagram for controller
with state limmitation and system}{fig:controller-fix}

Since neither the true gain of the system, or the motor timeconstant are known
precisely, they must be estimated as $G_e$ and $\tau_{motor,est}$. However as long
as the estimated values are close to the real values, and $\tau_{motor,est} \geq
\tau_{moto}$ and $G_e \leq G$ - that is the system is not estimated to be
faster, than it acually is - the system is still stabile. %jeg har plots og
% silieringer der verificere dette, men det fylder.  det kan også vises
% analytisk%

The only tuning needed is changing $\omega$. $\omega$ must be low so enough that
the phase shift introduced by the sensor filter is not an issue and low enough
for the output limmitation not to cause trouble. A value between $10$ and $25$
seemed good for our quadrotor.

\section{ Results } %change title?

The quadrotor was tethered to the legs of an up-side-down table, such that rotation
was restricted to the $y$-axis, for testing the attitude controller.

It became apparent during testing, that the mechanical vibration was so
powerful, that the IMU sensors saturate, resulting in instability due to bad
state estimation. In figure \ref{fig:omegaNoDamp8k} the gyro measurements
around the $y$-axis is plotted. It is very apparent that the measurements clip
at $\pm2^15$ when thrust is increased at about 17s.

\stdfignoscale{omegaNoDamp8k}{Gyro measuremenst around
$y$-axis without mechanical filter}{fig:omegaNoDamp8k}

To fix this, a quick mechanical filter made of foam earbuds was implemented.

In \ref{fig:omegaDamp8k} the raw gyro measurements after adding the mechanical filter
can be seen. At about 90s the throttle is raised to the same level, that caused
clipping before.

\stdfignoscale{omegaDamp8k}{Gyro measuremenst around
$y$-axis with mechanical filter}{fig:omegaDamp8k}

While this did solve the vibration problem, it also created some problems.
Sensor allignment becomes difficult, as the wires to the flight controller pulls the
flight controller ascew. Furthermore the mechanical filter has not been
modelled, so it's effect on the system is not accouted for.
Angular acceleration also seemes to suffer from this filter, in figure
\ref{fig:alphaDamp} and \ref{fig:alphaNoDamp} the angular accelleration
estimate, before the thrust is raised is clearly less noisy for the drone
without the mechanical filter.

\stdfignoscale{alphaDamp}{Angular acceleration around the
$y$-axis with mechanical filter}{fig:alphaDamp}
\stdfignoscale{alphaNoDamp}{Angular acceleration around the
$y$-axis with mechanical filter}{fig:alphaNoDamp}

However the angular speed measurements clearly benefit from this filter as seen
in \ref{fig:omegaDamp} and \ref{fig:omegaNoDamp}

\stdfignoscale{omegaDamp}{Angular speed around the
$y$-axis with mechanical filter}{fig:omegaDamp}
\stdfignoscale{omegaNoDamp}{Angular speed around the
$y$-axis with mechanical filter}{fig:omegaNoDamp}

A fix for this could be better balancing of propellers, as cheap propellers are
often slightly ansymmetrical. Another solution would be to find a more
predictable alternative for the earbuds.

In figure \ref{fig:thetaDamp} we see the estimated rotation around the
$y$-axis for the test with the mechanical filter. Zooming
(\ref{fig:thetaDampZoom}) in we can get a rough estimate of the bandwith of the
attitude controller. The input in this region is $[0^o, 45^o, -45^o, 0^o, 45^o,
0^o]$ and the rise time for the $45^o$ steps are roughly $1/3 s$ while the rise
time for the $90^o$ step is about $1/2 s$.

\stdfignoscale{thetaDamp}{Rotation around the
$y$-axis with mechanical filter}{fig:thetaDamp}
\stdfignoscale{thetaNoDamp}{Rotation around the
$y$-axis with mechanical filter, zoomed in on area of
interest.}{fig:thetaNoDamp}


A video of the drone in action can be found at: \url{youtu.be/jheuLw_xKc8}

\bibliographystyle{IEEEtran}
\bibliography{library}

%Novelties:
%*High sample-rate of acc/gyro, downsampled to remove aliasing.
% This results in propper attitude estimation, while allowing for high amounts
% of vibration.
%
% Mechanical dapening still needed to avoid sensor saturation.
%
% - plot of gyro before and after filter
% - plot of gyro with and widout earbuds

%*Hiding dynamics of rotor-spinup, resulting in easy controller design.
% - Copter simulink model
% - Bode plot w/wo controller (inner loop)

%*Controller designed by poleplacement, with state limitations.
% - Copter simulink model
% - Bode plot w/wo controller

%Camerastuff?


\end{document}
